\section{Silver and Gold Dataset}
\label{sec:silverandgold}

% TODO (~1.5--2 pages)
%
% --- 4.1 Silver labeling pipeline ---
% Tier structure (table: tier-0/1/2, signal, size, confidence):
%   Tier 0: 8,516,781 (92.44\%) — no ISBD signal → NER fallback; excluded from silver
%   Tier 1: 689,453 (7.48\%) — ISBD " : " delimiter present → high-confidence silver
%   Tier 2: ~7K — full ISBD structure (multiple delimiters) → highest confidence
% ISBD extraction rules (Stage 5):
%   " : " → splits TITLE from OTHER\_TITLE
%   " / " → splits title block from statement of responsibility → PERSON candidates
%   Role disambiguation: TRANSLATOR vs.\ PERSON (brief; cite sr04)
%   Trailing period noise: handling (cite sr05)
% Output: silver\_spans.jsonl — char-offset triples, same format as GLiNER training input.
% Exclusion: all 395 gold record IDs excluded from training data.
%
% --- 4.2 Gold annotation set ---
% Design: 395 records, stratified by era × silver tier (table: annotator × completed).
%   frankp: 393/395 (primary); maria: 131/395; gundi: 82/395 (IAA subsets)
%   ann (Claude pre-annotation, seeding only): not used for IAA.
% Annotation tool: Doccano; export format: JSONL char-offset triples.
% IAA: human-human pairs only — frankp∩maria (131 records), frankp∩gundi (81 records).
%   Treat one annotator as gold, other as predicted; span-level F1 + Cohen's κ.
%   Report range: [low, high] across both pairs.
% Reconciliation: frankp + user pass → gold\_reconciled.jsonl (due 17 June).
% Release: CC BY 4.0; silver spans + gold set + eval scripts at anonymous.4open.science/XXX.

\subsection{Gold Annotation Set}
\label{sec:gold}

A gold standard of 395 records was compiled from the DDB library sector,
stratified by publication era and silver labeling tier to ensure coverage
of the full range of orthographic and structural variation present in the
corpus (Table~\ref{tab:gold-strata}).
All 395 records were excluded from silver-label training data.

Three annotators with library cataloging expertise independently labeled
each record in Doccano, producing character-offset span triples in JSONL
format consistent with the silver output described in
Section~\ref{sec:silver}.
The label set is restricted to the three Work-level entities defined in
Section~\ref{sec:corpus}: \texttt{TITLE}, \texttt{OTHER\_TITLE}, and
\texttt{PERSON}.
Annotation coverage was near-complete across all three annotators
(Table~\ref{tab:gold-annotators}); records for which no labelable span
was found were explicitly marked as negative judgments and retained in
the data.

\begin{table}[h]
  \caption{Annotation coverage per annotator.}
  \label{tab:gold-annotators}
  \begin{tabular}{lrr}
    \toprule
    Annotator & Reviewed & of 395 \\
    \midrule
    A1 & 394 & 99.7\% \\
    A2 & 395 & 100\%  \\
    A3 & 392 & 99.2\% \\
    \bottomrule
  \end{tabular}
\end{table}

\paragraph{Inter-coder agreement.}
Inter-coder agreement was measured using Krippendorff's
$\alpha$~\cite{krippendorff1980,krippendorff2004a} with the nominal
distance function ($d_{ab} = 0$ if $a = b$, else $d_{ab} = 1$),
following the recommendations of \citet{artstein-poesio-2008}.
$\alpha$ was chosen for three reasons: it is defined for any number of
coders, accommodating the three-annotator setup directly; the nominal
distance function is the correct and unambiguous choice for mutually
exclusive entity labels; and, as a single-distribution coefficient,
$\alpha$ measures the reliability of the annotation \emph{procedure}
rather than the behaviour of individual annotators, which is the
appropriate target for a corpus released as a reusable
resource~\cite{artstein-poesio-2008}.
Agreement was computed on the pre-adjudication annotations to ensure the
coefficient reflects independent labeling.
The resulting $\alpha = \textbf{[TBD]}$ indicates \textbf{[TBD]}.

\paragraph{Reconciliation.}
Disagreements were resolved by a single pass of manual adjudication,
producing \texttt{gold\_reconciled.jsonl}, which is used in all
evaluations reported in Section~\ref{sec:results}.
The gold set, silver spans, and evaluation scripts are released under
CC~BY~4.0.\footnote{\url{https://anonymous.4open.science/XXX}}
